% Seccion: Objetivo e introducción
% Responsable: R1. Extension: 0,5 pagina(s).
% Solo R1 edita este archivo.
%
% Que debe contener (segun el enunciado y la pauta):
%   - Objetivo del laboratorio y del caso C.
%   - Organización del repositorio y cómo reproducir (docker compose run --rm lab1).
%   - Funciones auxiliares de bibliotecas usadas en el caso aplicado (np.fft, skimage: lectura y SSIM).

\section{Objetivo e introducción}
\label{sec:introduccion}

El laboratorio tiene dos objetivos. El primero es implementar desde cero, usando solo NumPy, las
herramientas de las Unidades 1--2 (muestreo, DFT, convolución y filtros FIR por ventana) y validarlas con
señales sintéticas cuyo resultado se conoce de antemano por la teoría. El segundo, el caso C, es usarlas
para estudiar Fourier en 2D: la DFT 2D y el papel de su magnitud y su fase, la eliminación de interferencia
periódica con filtros notch (con PSNR y SSIM contra la imagen original, y sobre fotos reales sin referencia)
y el aliasing espacial al reducir una imagen sin filtrarla antes.

\paragraph{Organización y reproducción.}
{\emergencystretch=3em  % evita que un nombre largo en \texttt se salga del margen
El paquete \texttt{psi\_lab.\allowbreak lab1} tiene un módulo por herramienta (\texttt{muestreo}, \texttt{dft},
\texttt{convolucion}, \texttt{filtros}, \texttt{ventanas}, \texttt{sinteticos}), un subpaquete
\texttt{caso\_imagen} para el caso C y un archivo por figura en \texttt{figuras/}. Cada figura y tabla de
este informe sale del código: \texttt{docker compose run --rm lab1} regenera todas en \texttt{out/}, y
descarga las dos fotos reales del caso C desde un \emph{release} del repositorio (no se versionan datos).
\texttt{docker compose run --rm tests} corre las pruebas unitarias (propiedades teóricas, no valores
copiados) y de integración; el CI ejecuta además el análisis estático y las pruebas en Python 3.11 a 3.13.\par}

\paragraph{Funciones de bibliotecas usadas.}
En la parte común solo se usa NumPy para el cálculo; \texttt{np.fft}, \texttt{np.convolve} y
\texttt{scipy} aparecen solo en las pruebas, como referencia (y \texttt{np.fft.fft} en la comparación
de tiempos de la DFT). En el caso C usamos como auxiliares \texttt{np.fft} para la DFT 2D y
\texttt{skimage} para cargar imágenes y calcular SSIM. La detección de picos, los notch, el PSNR y
el pasabajos gaussiano son propios.
